% from ega4-16-fr.tex
\oldpage[IV-4]{47}

Les deux conditions de l'énoncé ne font que traduire la définition donnée dans
(0, 15.2.2), compte tenu de la définition des homomorphismes canoniques
(0, 15.2.1.1). La dernière assertion résulte de ce que, si un module $M$ sur
un anneau commutatif $A$ admet une base de $n$ éléments, tout système de
générateurs de $M$ ayant $n$ éléments est une base de $M$ (Bourbaki,
\emph{Alg. comm.}, chap.~II, §~3, cor.~5 du th.~1).

\begin{corollary}[16.9.4]
\label{IV.16.9.4}
Soient $X$ un espace annelé en anneaux locaux, $\mathscr{I}$ un Idéal de
$\mathcal{O}_X$. Pour que $\mathscr{I}$ soit quasi-régulier, il faut et il
suffit qu'il vérifie les conditions suivantes:
\begin{enumerate}
  \item[(i)] $\mathscr{I}$ est de type fini.
  \item[(ii)] $\mathscr{I}/\mathscr{I}^2$ est un
  $(\mathcal{O}_X/\mathscr{I})$-Module localement libre.
  \item[(iii)] L'homomorphisme canonique
  \[
    \mathbf{S}_{\mathcal{O}_X/\mathscr{I}}^{\bullet}
    (\mathscr{I}/\mathscr{I}^2)
    \longrightarrow\mathscr{G}r_{\mathscr{I}}^{\bullet}(\mathcal{O}_X)
    \tag{16.9.4.1}
    \label{IV.16.9.4.1}
  \]
  est bijectif.
\end{enumerate}
\end{corollary}

La nécessité des conditions résulte aussitôt de (16.9.3). Pour voir que ces
conditions sont suffisantes, il faut montrer, en vertu de (16.9.3), que si, en
un point $x\in\operatorname{Supp}(\mathcal{O}_X/\mathscr{I})$, il existe un
voisinage ouvert $U$ de $x$ dans $X$ et $n$ sections $f_i$ ($1\leq i\leq n$)
de $\mathscr{I}$ au-dessus de $U$ dont les images canoniques dans
$\mathscr{I}/\mathscr{I}^2$ forment une base de
$(\mathscr{I}/\mathscr{I}^2)|U$ sur
$(\mathcal{O}_X/\mathscr{I})|U$, alors il y a un voisinage ouvert
$V\subset U$ de $x$ tel que les $f_i|V$ engendrent $\mathscr{I}|V$. Or, par
hypothèse, on a $\mathscr{I}_x\neq\mathcal{O}_x$, donc $\mathscr{I}_x$ est
contenu dans l'idéal maximal de $\mathcal{O}_x$; comme $\mathscr{I}_x$ est un
$\mathcal{O}_x$-module de type fini et que les classes des $(f_i)_x$ dans
$\mathscr{I}_x/\mathscr{I}_x^2$ engendrent cet
$(\mathcal{O}_x/\mathscr{I}_x)$-module, le lemme de Nakayama montre que les
$(f_i)_x$ engendrent $\mathscr{I}_x$. Comme $\mathscr{I}$ est de type fini,
on conclut par (0\textsubscript{I}, 5.2.2).

\begin{corollary}[16.9.5]
\label{IV.16.9.5}
Soient $X$ un espace annelé en anneaux locaux, $\mathscr{I}$ un Idéal
quasi-régulier de $\mathcal{O}_X$, $(f_i)_{1\leq i\leq n}$ une suite de
sections de $\mathscr{I}$ au-dessus de $X$,
$x$ un point de $\operatorname{Supp}(\mathcal{O}_X/\mathscr{I})$. Les
conditions suivantes sont équivalentes:
\begin{enumerate}
  \item[(a)] Il existe un voisinage ouvert $U$ de $x$ dans $X$ tel que les
  $f_i|U$ forment une suite quasi-régulière d'éléments de
  $\Gamma(U,\mathcal{O}_X)$ engendrant $\mathscr{I}|U$.
  \item[(b)] Les $(f_i)_x$ forment un système de générateurs de
  $\mathscr{I}_x$ dont le nombre d'éléments est le plus petit possible.
  \item[(b')] Les $(f_i)_x$ forment un système minimal de générateurs de
  $\mathscr{I}_x$.
  \item[(c)] Si $\bar f_i$ est l'image canonique de $f_i$ dans
  $\Gamma(X,\mathscr{I}/\mathscr{I}^2)$, les $(\bar f_i)_x$ forment une base
  du $(\mathcal{O}_x/\mathscr{I}_x)$-module
  $\mathscr{I}_x/\mathscr{I}_x^2$.
\end{enumerate}
\end{corollary}

Par hypothèse, $\mathcal{O}_x$ est un anneau local, $\mathscr{I}_x$ un idéal
de type fini de $\mathcal{O}_x$ contenu dans l'idéal maximal de
$\mathcal{O}_x$; l'équivalence de b), b') et c) résulte donc du lemme de
Nakayama (Bourbaki, \emph{Alg. comm.}, chap.~II, §~3,
n\textsuperscript{o}~2, prop.~5). Il est clair que a) entraîne c) en vertu de
(16.9.3); d'autre part, il résulte de (0\textsubscript{I}, 5.2.2) que si la
condition c) est vérifiée (donc aussi b)), il y a un voisinage $U$ de $x$ dans
$X$ tel que $(\mathscr{I}/\mathscr{I}^2)|U$ ait un rang constant égal à $n$,
et que les $f_i|U$ engendrent $\mathscr{I}|U$; il suffit donc d'appliquer dans
$U$ la dernière assertion de (16.9.3).

\begin{remarks}[16.9.6]
\label{IV.16.9.6}
\begin{enumerate}
  \item[(i)] Sous les hypothèses générales de (16.9.5), il ne suffit pas que
  les $(\bar f_i)_y$ forment une base du
  $(\mathcal{O}_y/\mathscr{I}_y)$-module
  $\mathscr{I}_y/\mathscr{I}_y^2$ pour tout $y\in X$ pour que la suite
  $(f_i)$ engendre $\mathscr{I}$. On en a un exemple en
